\section{Conclusiones}
\label{sec:conclusiones}

\subsection{Principales resultados}
\begin{itemize}
    \item \underline{AM:} se identificaron Cadena 3 ($700 \, [kHz]$) y Radio Mitre ($810 \, [kHz]$), con portadoras de unos $-101$ y $-102 \, [dBm]$, apenas $3 \, [dB]$ por encima del piso de ruido. Con la antena de varillas, pensada para VHF, las bandas laterales quedaron ocultas, por lo que el ancho de banda ocupado y la máxima frecuencia modulante sólo se analizaron en forma teórica (como máximo $10 \, [kHz]$ y unos $5 \, [kHz]$, según la canalización). Con el generador, la modulante cuadrada produjo el peine de bandas laterales en las armónicas impares, separadas $2 f_m = 20 \, [kHz]$, que predice su serie de Fourier; con la senoidal, el RBW de $9 \, [kHz]$ no permitió resolver las bandas laterales a $10 \, [kHz]$ de la portadora, y el índice de modulación no pudo verificarse.
    \item \underline{FM:} se distinguieron alrededor de una docena de emisoras, con niveles de $-90$ a $-65 \, [dBm]$. La emisora de $88{,}9 \, [MHz]$ ocupó unos $220 \, [kHz]$ con el criterio de $-26 \, [dB]$: algo más que la separación de $200 \, [kHz]$ entre canales y entre los valores de Carson para una emisión monofónica ($180 \, [kHz]$) y una estéreo ($256 \, [kHz]$). Junto a ella se observó una emisora en el canal adyacente, con los flancos de ambas casi en contacto cerca del piso de ruido. La captura en barrido simple de la emisora de $102{,}3 \, [MHz]$ mostró que el espectro instantáneo de una FM con programa cambia en cada barrido.
    \item \underline{Televisión analógica:} se recibieron los canales 8, 10 y 12, con los canales intermedios 9 y 11 libres, como prevé la planificación. En los canales 10 y 12 se ubicaron la portadora de video, a $1{,}25 \, [MHz]$ del borde inferior, y la de audio, $4{,}5 \, [MHz]$ más arriba y unos $10 \, [dB]$ por debajo; en el canal 12 se distinguió además la subportadora de color. Las bandas laterales de video quedaron bajo el piso de ruido, por lo que su ancho de banda no pudo medirse.
    \item \underline{TDA:} se observó el canal 27 como una meseta de unos $5{,}5 \, [MHz]$, coherente con los $5{,}57 \, [MHz]$ de la señal ISDB-Tb, sin portadora destacada y con una potencia de canal estimada en unos $-61 \, [dBm]$, unos $11 \, [dB]$ por encima del piso de ruido del analizador. No se realizó la prueba con un televisor, por lo que no se compararon las señales observadas con los servicios decodificados.
    \item \underline{Niveles:} la conversión entre dBm y dBµV depende de la impedancia ($+107{,}0 \, [dB]$ en $50 \, [\Omega]$ y $+108{,}8 \, [dB]$ en $75 \, [\Omega]$), mientras que los niveles en dBc no dependen de ella. En el ejercicio, una banda lateral de $-40 \, [dBc]$ corresponde a un índice de modulación del $2 \,\text{\%}$.
    \item \underline{Distribución por cable y fibra:}
    \begin{itemize}
        \item DVB-C aprovecha la alta relación señal a ruido del cable con QAM de portadora única, y transporta unos $38 \, [Mbps]$ por canal de $6 \, [MHz]$ con 256-QAM.
        \item DOCSIS convierte la red HFC en una red de datos bidireccional, y es la base del acceso a Internet y de la migración del video a IP.
        \item El RF Overlay lleva todo el espectro de una red de cable en $1.550 \, [nm]$ sobre la ODN de una red FTTH, sin consumir capacidad de datos y sin decodificadores, a costa de equipamiento óptico de video y sin interactividad. La IPTV, en cambio, integra el video en la red de datos.
        \item ISDB-Tb usa OFDM en 13 segmentos y capas jerárquicas, que permiten combinar en un canal servicios fijos en HD y servicios móviles.
        \item ISDB-Tb sobre cable aprovecha que los televisores argentinos ya traen sintonizador ISDB-Tb para digitalizar las redes de cable sin decodificadores, a cambio de una capacidad por canal menor que la de DVB-C.
    \end{itemize}
\end{itemize}

\subsection{Dificultades y limitaciones de las mediciones}
\begin{itemize}
    \item \underline{Antena:} la antena interior de varillas telescópicas es adecuada para VHF, pero en la banda de AM es eléctricamente muy corta, lo que dejó a las emisoras de AM apenas por encima del piso de ruido.
    \item \underline{Piso de ruido del analizador:} en torno a $-100 \, [dBm]$, ocultó las bandas laterales de AM, las de video de la televisión analógica y, posiblemente, otros canales digitales de bajo nivel.
    \item \underline{Resolución del instrumento:} el RBW de $9 \, [kHz]$ es comparable con el canal de AM y con la separación de las bandas laterales de la señal generada, y no las resuelve; un RBW menor hace más lento el barrido.
    \item \underline{Señales de programa:} su espectro varía todo el tiempo, por lo que los anchos de banda medidos dependen del tiempo de observación, de la retención de máximos y del criterio de nivel adoptado.
    \item \underline{Lectura de los valores:} los niveles y frecuencias se leyeron sobre la grilla de fotografías de la pantalla, por lo que son aproximados; con los marcadores del analizador se habría obtenido mayor precisión.
    \item \underline{Impedancia:} el analizador mide en $50 \, [\Omega]$ y las antenas y redes de televisión son de $75 \, [\Omega]$. La desadaptación o el adaptador agregan un error o una atenuación que hay que considerar.
    \item \underline{Equipamiento:} no se dispuso de un televisor o receptor ISDB-Tb, por lo que no se pudo comprobar la recepción de la TDA.
\end{itemize}
